\documentclass[12pt]{article}

\usepackage[paperwidth=612pt,paperheight=792pt,top=72pt,bottom=72pt,left=30mm,right=20mm]{geometry}
\usepackage{times}
\usepackage{setspace}
\usepackage{fancyhdr}
\usepackage{ragged2e}
\usepackage{graphicx}
\usepackage{mdframed}
\usepackage{enumitem}
\usepackage{listings}
\setstretch{1.5}

\pagestyle{fancy}
\fancyhf{}
\fancyhead[R]{AyurSpace}
\fancyfoot[L]{Pillai HOC College of Engineering \& Technology (Autonomous), B. E. Computer Engineering}
\fancyfoot[R]{\thepage}
\renewcommand{\headrulewidth}{1pt}
\renewcommand{\footrulewidth}{1pt}

\begin{document}

\vspace*{\fill}
\begin{center}
\textbf{\Huge{Chapter 4}}\\
\vspace{5mm}
\textbf{\Huge{Plan Of Project}}
\end{center}
\vspace*{\fill}
\newpage

\section*{4.1 \quad Methodology}
\addcontentsline{toc}{section}{4.1 Methodology}
\justify

AyurSpace uses a highly decoupled method that combines artificial intelligence, mobile frontend
engineering, and backend cloud infrastructure.

\begin{itemize}
    \item \textbf{Presentation Layer (Flutter):}\\
    Built entirely in Flutter, responsible for UI rendering, capturing user input, translating
    languages (ARB files), and displaying state updates effectively across Android and iOS platforms.

    \item \textbf{State Management:}\\
    We utilize Riverpod for reactive state management and dependency injection. Riverpod isolates UI
    from internal business logic, enabling components to listen to streams for live changes (e.g.,
    chat message history, dosha quiz progression).

    \item \textbf{Domain Layer (Core Business Logic):}\\
    Contains pure Dart classes mapping Ayurvedic concepts into technical objects---like
    \texttt{Plant}, \texttt{Dosha}, and \texttt{Remedy}.

    \item \textbf{Plant Taxonomy Integration (Plant.id API):}\\
    When a user scans a plant, the photo is compressed locally and sent to the Plant.id residual
    network, acting as a taxonomic classifier identifying the technical species name.

    \item \textbf{LLM Contextualizing (Gemini API):}\\
    After technical identification, a prompt is generated for Google Gemini API, employing
    Chain-of-Thought (CoT) prompting techniques. Gemini receives the Latin name to synthesize
    Ayurvedic properties such as \textit{Rasa}, \textit{Guna}, \textit{Virya}, and home remedies.

    \item \textbf{Cloud Storage and Authentication:}\\
    To ensure accurate user profiling across devices, we employ Firebase Authentication and Cloud
    Firestore. It powers live synchronization of users' scanning history, determined dosha
    conditions, and profile data securely.
\end{itemize}

\newpage

\section*{4.2 \quad Project Plan (Timeline Diagram)}
\addcontentsline{toc}{section}{4.2 Project Plan}

A Gantt chart illustrates the timeline of various project modules from Planning to Deployment.

\renewcommand{\thefigure}{4.1}
\renewcommand{\figurename}{Figure}
\begin{figure}[!htbp]
\begin{center}
\includegraphics[width=\textwidth,height=0.4\textheight,keepaspectratio]{GanttChart.png}
\caption{Project Plan Timeline}
\end{center}
\end{figure}

The timeline ensures efficient project management, tracking module initiation and deadlines for the
UI layout, API integration, testing, and deployment.

\newpage

\section*{4.3 \quad Proposed System Architecture}
\addcontentsline{toc}{section}{4.3 Proposed System Architecture}
\justify

The AyurSpace system architecture is built as a cloud-native mobile application that follows Clean
Architecture principles. The system works on a Client-Server model, where the Flutter mobile client
is the orchestration engine that links different microservices. The application logic is highly
decoupled from the presentation layer, relying on dependency injection and robust state management.

\subsection*{4.3.1 \quad Layered Architecture}
The software consists of three primary layers:
\begin{enumerate}
    \item \textbf{Presentation Layer}: Built entirely in Flutter. Responsible for UI rendering,
    capturing user input, and displaying state updates.

    \item \textbf{Domain Layer}: Contains the core business logic, entities, and use-case
    definitions. Functions independently of any specific data source or UI framework.

    \item \textbf{Data Layer}: Handles the retrieval and storage of data from local caches,
    Firebase, and third-party APIs like Plant.id and Google Gemini.
\end{enumerate}

\subsection*{4.3.2 \quad Directory Structure}
The repository structure enforces this separation of concerns:
\begin{verbatim}
lib/
|-- config/          # Global configs (API keys, theme, router)
|-- data/            # Models, Repositories, Data Sources
|   |-- models/
|   |-- repositories/
|   |-- sources/
|-- providers/       # Riverpod Providers and Notifiers
|-- screens/         # UI Screens
|-- services/        # Service wrappers (API clients, Firebase)
|-- utils/           # Utilities and helpers
|-- widgets/         # Reusable UI components
\end{verbatim}

\subsection*{4.3.3 \quad Hybrid Neuro-Symbolic Inference Pipeline}
The system's intelligence relies on a two-step hybrid approach:
\begin{itemize}
    \item \textbf{Discriminative Phase (Vision)}: An image uploaded by the user is compressed and
    sent to the Plant.id API. The API uses a CNN to return the scientific classification (Taxonomy)
    with a confidence score.

    \item \textbf{Generative Phase (Language)}: The scientific name identified by the CNN is passed
    as a prompt to the Google Gemini LLM. Gemini is constrained via strict system prompts to act as
    an Ayurvedic expert, returning the \textit{Dosha} properties, uses, and precautions associated
    with the plant.
\end{itemize}

\renewcommand{\thefigure}{4.2}
\begin{figure}[!htbp]
\begin{mdframed}
    \centering
    \includegraphics[width=0.7\linewidth,keepaspectratio]{SystemArch.png}
\end{mdframed}
\caption{Proposed System Architecture}
\end{figure}

\end{document}
